\documentclass[11pt]{article}
\usepackage[utf8]{inputenc}
\usepackage[T1]{fontenc}
\usepackage[spanish]{babel}
\usepackage{amsmath,amssymb}
\usepackage{geometry}
\usepackage{booktabs}
\usepackage{hyperref}
\geometry{letterpaper, margin=1in}

\title{Modelo local de red narrativa ponderada y cubridor nodal}
\author{}
\date{}

\begin{document}
\maketitle

\section{Objeto y alcance}

Se propone un sistema local para analizar narrativas en corpus documentales.
El sistema no envía textos a modelos externos. Extrae documentos, limpia y
normaliza texto, clasifica fuentes, identifica estructura narrativa y construye
una red compleja ponderada para resolver un cubridor nodal narrativo. El objetivo
no es inferir ``la narrativa total'' de un tema, sino producir una muestra
auditable, filtrable y reproducible de relaciones narrativas.

El modelo adopta una lógica de \textit{Smart Data}: antes de optimizar, separa
capas de fuente, elimina ruido, identifica marcos argumentales, detecta ecos
documentales, registra cambios temporales y conserva trazabilidad técnica. Por
lo tanto, el cubridor no se aplica sobre texto crudo, sino sobre una red ya
normalizada, estratificada y revisable.

La tubería separa explícitamente cuatro géneros de fuente: artículos
científicos abiertos, noticias, foros/discusiones públicas y fuentes
institucionales o industriales. Esta separación evita mezclar evidencia de
distinta naturaleza como si tuviera el mismo peso.

\section{Preprocesamiento y extracción narrativa}

Antes del análisis lingüístico, los textos se homogeneizan en minúsculas, se
normalizan acentos, se eliminan URL, puntuación no informativa y palabras vacías
en español e inglés. La extracción de monogramas, bigramas, trigramas y redes de
coocurrencia se realiza sobre este texto normalizado. Para detectar personas u
organizaciones se conserva también el texto editorial original, pero se filtran
etiquetas metodológicas y ruido como \textit{abstract}, \textit{conclusion},
\textit{journal}, nombres de fuente y términos que son el propio tópico de
búsqueda.

Para cada documento $d_i$ se extraen cinco etapas narrativas:
\[
E=\{\text{evento inicial},\text{conflicto},\text{punto de cambio},
\text{resolución},\text{consecuencias}\}.
\]
También se extraen actores, fuente, año, localización y tipo de fuente. La
extracción es heurística y auditable: cada etapa conserva la oración detectada y
los marcadores textuales que activaron la clasificación. Los actores se tratan
como candidatos, no como verdad automática.

Además, se extraen proposiciones simples de la forma sujeto--verbo--objeto,
actos de habla aproximados, marcadores de causalidad y señales argumentales. Las
salidas no se interpretan como prueba de verdad, sino como una lista de lugares
del corpus que requieren lectura humana.

\section{Núcleo Smart Data y ontología de relaciones}

Antes de construir el cubridor, el sistema organiza la información en cuatro
capas: cartografía de fuentes, marcos argumentales, evolución temporal y
silencios. Para cada documento se estiman marcos del tipo:
\[
M_d=(problema_d, culpable_d, solucion_d, urgencia_d),
\]
cuando el texto contiene señales suficientes. Si varios documentos repiten el
mismo marco, se registra como eco y no como diversidad narrativa adicional.

La red distingue relaciones como:
\[
{\mathcal R}=\{
\text{afirma},\text{apoya},\text{critica},\text{causa},\text{amplifica},
\text{contradice},\text{silencia},\text{contextualiza}
\}.
\]
Estas etiquetas son operativas y revisables. No prueban intencionalidad; sirven
para ordenar el análisis y hacer visible qué tipo de conexión sostiene cada
arista.

\section{Validación humana de actores}

Los actores automáticos no se consideran definitivos. El sistema incorpora una
intervención humana para aceptar, corregir, rechazar o agregar actores por
documento. La validación se guarda localmente en \texttt{actor\_validation.json}
con estado, validador, notas y fecha. Cuando existe validación, la red narrativa
usa los actores validados; si el registro es rechazado, el documento no aporta
actores a la red.

\section{Red narrativa}

La red narrativa se define como:
\[
G=(V,A,w,\tau,\sigma),
\]
donde $V$ es el conjunto de nodos, $A$ el conjunto de aristas ponderadas,
$w:A\to \mathbb{R}^{+}$ asigna peso a cada relación, $\tau:V\to{\mathcal T}$
indica el tipo de nodo y $\sigma:V\to \mathbb{R}^{+}$ asigna importancia
estructural-narrativa al nodo.

Los nodos incluyen:
\[
V=V_D\cup V_R\cup V_Q\cup V_F\cup V_T\cup V_L\cup V_C\cup V_N,
\]
donde $V_D$ son documentos, $V_R$ actores, $V_Q$ etapas narrativas, $V_F$
fuentes, $V_T$ años, $V_L$ localizaciones, $V_C$ tipos de fuente y $V_N$
conceptos lingüísticos: monogramas, bigramas y trigramas.

Además se define una capa de flujo narrativo:
\[
F=\{\text{diálogo individual},\text{noticia},\text{derivación técnica},
\text{investigación},\text{otras derivaciones}\}.
\]
Esta capa no impone peso automático. Sirve para observar circulación: el diálogo
individual puede preceder a la noticia, la noticia puede preceder a la
investigación, y los resultados de investigación pueden retroalimentar diálogos
y noticias posteriores.

Las aristas representan relaciones como:
\begin{itemize}
\item documento--actor, documento--etapa y documento--concepto;
\item actor--etapa, fuente--etapa, año--etapa y localización--etapa;
\item fuente--concepto y coocurrencia concepto--concepto.
\end{itemize}
Cada arista $a\in A$ se representa como una terna:
\[
a=(i,j,\rho),\quad i,j\in V,\quad \rho\in{\mathcal R},
\]
donde $\rho$ es el tipo de relación. Para evitar introducir una jerarquía
artificial entre elementos de narrativa social, el modo base usa ponderación
neutral:
\[
w_a = \sum_{d\in D}\mathbf{1}\{a \text{ aparece en } d\}.
\]
Así, una relación pesa más sólo porque aparece más veces en el corpus filtrado,
no porque el investigador haya decidido de antemano que una etapa, fuente o
actor vale más que otro. Los tipos de fuente se usan para estratificar y
comparar resultados. Si se usa un peso de evidencia \(\omega_d\), debe
declararse como análisis de sensibilidad:
\[
w_a^{sens} = \sum_{d\in D}\mathbf{1}\{a \text{ aparece en } d\}\cdot \omega_d.
\]
Esta distinción es importante porque distintos pesos producen sesgos distintos
en la red y, por tanto, en el cubridor.

Para comparación estructural se conserva el conteo crudo y se añade una versión
normalizada:
\[
\widetilde{w}_a=\frac{w_a}{\max_{b\in A}w_b}\in[0,1],
\qquad
\widetilde{c}_v=\frac{c_v}{\max_{u\in V}c_u}\in[0,1],
\]
donde \(c_v\) es el conteo de apariciones del nodo. El grado ponderado también
se normaliza:
\[
\widetilde{d}^{\,w}_v=
\frac{\sum_{a\sim v}w_a}{\max_{u\in V}\sum_{b\sim u}w_b}.
\]
Así, los pesos usados para comparar redes oscilan entre 0 y 1, mientras que los
conteos crudos quedan disponibles para auditoría.

Antes de crear los grafos semántico y de conocimiento se realiza composición de
frases. Si un monograma o bigrama está explicado mayoritariamente por una frase
repetida más larga, se marca como absorbido:
\[
q \prec p
\quad \text{si}\quad
q \subset p
\ \text{y}\ 
\frac{freq(p)}{freq(q)}\geq \theta.
\]
Por ejemplo, si \textit{violencia social} explica la mayoría de apariciones de
\textit{violencia} y \textit{social}, el nodo principal es la frase completa y
las aristas de coocurrencia y documento--concepto se construyen con esa unidad
canónica. Para evitar conceptos accidentales, la frase canónica debe repetirse
al menos dos veces en el corpus filtrado.

La importancia nodal usada por el cubridor combina frecuencia y estructura, sin
ponderar por tipo narrativo en el escenario base:
\[
\sigma_v =
\gamma_1\widetilde{c}_v+
\gamma_2\widetilde{d}^{\,w}_v+
\gamma_3\widehat{ev}_v+
\gamma_4\widehat{val}_v,
\]
donde $\widehat{freq}_v$ es frecuencia normalizada, $\widehat{deg}^{\,w}_v$ es
grado ponderado normalizado, $\widehat{ev}_v$ es evidencia promedio asociada al
nodo y $\widehat{val}_v$ registra validación humana. En el análisis base se
recomienda fijar \(\gamma_3=0\) si se mezclan géneros de fuente, y usar
\(\widehat{ev}_v\) sólo en corridas de sensibilidad o en análisis separado por
tipo de fuente. Los parámetros $\gamma_i\geq 0$ satisfacen $\sum_i\gamma_i=1$.

\section{Selector nodal multiobjetivo inspirado en cobertura}

El problema se formula como una adaptación inspirada en el \textit{set covering
problem} (SCP), pero no debe confundirse con el SCP clásico estricto si no se
exige cubrir todo el universo de relaciones. Sea el universo operativo que se
desea explicar:
\[
{\mathcal U}=A^\star\subseteq A,
\]
donde $A^\star$ se obtiene aplicando filtros metodológicos: tipos de arista
permitidos, peso mínimo, fuentes incluidas y exclusiones manuales. Cada nodo
candidato $v\in B\subseteq V$ define un subconjunto:
\[
S_v=\{a\in A^\star: a \text{ es incidente a } v\}.
\]
La familia de subconjuntos es:
\[
{\mathcal S}=\{S_v:v\in B\}.
\]
Una cobertura clásica es una subfamilia:
\[
{\mathcal C}\subseteq{\mathcal S}
\]
tal que:
\[
\bigcup_{S_v\in{\mathcal C}}S_v={\mathcal U}.
\]
En términos nodales, seleccionar ${\mathcal C}$ equivale a seleccionar un
subconjunto de nodos $C\subseteq B$.

La versión de decisión pregunta si existe una cobertura con $|C|\leq k$. Esta
forma corresponde al problema clásico NP-completo de cobertura de conjuntos. En
la aplicación narrativa se usa una versión de optimización multiobjetivo: se
busca un conjunto de nodos pequeño e importante, pero se penaliza el daño
estructural producido si ese conjunto se retira de la red.

\subsection{Variables de decisión}

Para cada nodo candidato $v\in B$:
\[
x_v =
\begin{cases}
1, & \text{si el nodo }v\text{ es seleccionado},\\
0, & \text{en otro caso}.
\end{cases}
\]
Para cada arista objetivo $a\in A^\star$:
\[
r_a =
\begin{cases}
1, & \text{si la arista }a\text{ se elimina al retirar }C,\\
0, & \text{en otro caso}.
\end{cases}
\]
Definimos además la matriz de incidencia:
\[
b_{av} =
\begin{cases}
1, & \text{si }a\in S_v,\\
0, & \text{en otro caso}.
\end{cases}
\]

\subsection{Parámetros}

\begin{itemize}
  \item $n=|B|$: número de nodos candidatos.
  \item $m=|A^\star|$: número de aristas objetivo.
  \item $k$: máximo de nodos seleccionables.
  \item $w_a>0$: peso de la arista $a$.
  \item $\sigma_v>0$: importancia del nodo $v$.
  \item $e_v\in\{0,1\}$: elegibilidad del nodo; vale $0$ si el usuario excluye
  el nodo, el medio, el término o el tipo de fuente.
  \item $q_{vh}\in\{0,1\}$: pertenencia del nodo $v$ al tipo $h$.
  \item $L_h,U_h$: mínimos y máximos opcionales por tipo de nodo.
  \item $\eta\in[0,1]$: relevancia mínima ponderada opcional.
\end{itemize}

Sean:
\[
W=\sum_{a\in A^\star}w_a,\qquad
S=\sum_{v\in B}\sigma_v.
\]

\subsection{Restricciones}

Las restricciones binarias son:
\[
x_v\in\{0,1\}\quad \forall v\in B,\qquad
r_a\in\{0,1\}\quad \forall a\in A^\star.
\]

La remoción lógica de aristas se impone con:
\[
r_a \geq b_{av}x_v
\quad \forall a\in A^\star,\ \forall v\in B,
\]
\[
r_a \leq \sum_{v\in B} b_{av}x_v
\quad \forall a\in A^\star.
\]
Así, una arista se cuenta como removida si toca al menos un nodo seleccionado:
\[
A_R(C)=\{a=(i,j)\in A^\star:i\in C \lor j\in C\}.
\]

El presupuesto de complejidad interpretativa es:
\[
\sum_{v\in B}x_v \leq k.
\]

La exclusión manual y metodológica se modela como:
\[
x_v\leq e_v\quad \forall v\in B.
\]

Cuando se requiere controlar la composición del cubridor por tipo de nodo:
\[
L_h \leq \sum_{v\in B}q_{vh}x_v \leq U_h
\quad \forall h\in{\mathcal T}.
\]
Por ejemplo, puede limitarse el número de documentos para evitar soluciones
triviales y exigir presencia de actores, etapas o conceptos.

\subsection{Objetivos}

El modelo compara soluciones mediante tres objetivos normalizados:
\[
f_1(x)=\frac{\sum_{v\in B}x_v}{n}
\quad \text{minimizar proporción de nodos seleccionados},
\]
\[
f_2(x)=\frac{\sum_{v\in B}\sigma_v x_v}{S}
\quad \text{maximizar proporción de importancia nodal retenida},
\]
\[
f_3(r)=\frac{\sum_{a\in A^\star}w_a r_a}{W}
\quad \text{minimizar proporción de peso de aristas removidas}.
\]

El problema multiobjetivo queda:
\[
\min f_1(x),\qquad \max f_2(x),\qquad \min f_3(r),
\]
sujeto a las restricciones anteriores. Para trabajar en un espacio de
maximización común se usa:
\[
g(x,r)=\left(1-f_1(x),\ f_2(x),\ 1-f_3(r)\right)\in[0,1]^3.
\]

\subsection{Dominancia de Pareto e hipervolumen}

Una solución $(x^1,r^1)$ domina a $(x^2,r^2)$ si:
\[
g_i(x^1,r^1)\geq g_i(x^2,r^2)\quad \forall i
\]
y existe al menos un objetivo $j$ tal que:
\[
g_j(x^1,r^1)>g_j(x^2,r^2).
\]
El frente de Pareto $P$ contiene las soluciones factibles no dominadas.

Para comparar métodos se estima un hipervolumen normalizado:
\[
HV(P)=\lambda_3\left(\bigcup_{(x,r)\in P}
[r_1,g_1(x,r)]\times[r_2,g_2(x,r)]\times[r_3,g_3(x,r)]\right),
\]
donde $r=(0,0,0)$ es el punto de referencia y $\lambda_3$ es la medida de
volumen en tres dimensiones. En la implementación local este volumen se aproxima
con muestreo Monte Carlo determinístico. Un mayor valor indica un frente que
combina mejor parsimonia, importancia nodal y bajo daño estructural por remoción
de aristas.

\subsection{Función escalar auxiliar}

Los algoritmos heurísticos requieren ordenar candidatos durante la exploración.
Para ello se usa una función auxiliar:
\[
\Phi(x,r)=
\alpha f_2(x)
-\beta f_3(r)
-\lambda f_1(x)
-\sum_h \pi_h(x),
\]
donde $\alpha,\beta,\lambda\geq0$ y $\pi_h(x)$ penaliza violaciones o
desbalances respecto a la composición deseada por tipo de nodo. Esta función no
reemplaza el modelo multiobjetivo: sólo guía exploraciones auxiliares y barridos
para generar soluciones candidatas cuyo desempeño final se evalúa por Pareto e
hipervolumen.

\section{Métodos de solución}

El sistema implementa cuatro métodos locales. Para que la comparación sea
válida, todos se ejecutan sobre la misma instancia \((G,B,A^\star,k)\), con las
mismas restricciones, el mismo criterio de factibilidad, el mismo presupuesto
\(E\) de evaluaciones de función objetivo y las mismas métricas de calidad de
frente. Sólo cambia la estrategia de búsqueda:

\begin{enumerate}
  \item \textbf{Barrido glotón por suma ponderada}: recorre vectores de peso
  sobre \(u_1,u_2,u_3\) y construye soluciones incrementales.
  \item \textbf{MOEA}: mantiene una población de subconjuntos, cruza, muta y
  conserva un archivo de soluciones no dominadas.
  \item \textbf{MOSA}: usa recocido simulado multiobjetivo; acepta soluciones por
  dominancia y temperatura cuando la comparación no es concluyente.
  \item \textbf{MMC-MO}: adaptación discreta multiobjetivo del método de
  composición musical de Mora-Gutiérrez et al.; cada compositor conserva un
  arreglo de nodos y modifica su arreglo usando memoria guía LP-Pareto.
\end{enumerate}

\section{Adaptación del MMC}

En el MMC original, una sociedad de compositores produce arreglos musicales. En
este problema, cada arreglo es un subconjunto de nodos narrativos:
\[
A_c=\{v_1,\ldots,v_m\},\quad m\leq k.
\]
La calidad del arreglo se evalúa mediante el vector multiobjetivo y su aporte al
frente de Pareto. El genio $G^*$ es una solución guía seleccionada entre las
mejores soluciones no dominadas. En cada iteración se aplican tres operadores:

\begin{itemize}
  \item innovación guiada por genio: insertar nodos de $G^*$;
  \item cambio: agregar, eliminar o sustituir nodos;
  \item intercambio social: recombinar el arreglo de un compositor con otro.
\end{itemize}

La solución nueva reemplaza a la anterior usando el mismo criterio común de
dominancia con factibilidad:
entre dos soluciones factibles se aplica Pareto; entre una factible y una
infactible se conserva la factible; entre dos infactibles se prefiere la de
menor violación total de restricciones. El desempeño global de todos los métodos
se reporta con las mismas métricas:
\begin{itemize}
  \item hipervolumen normalizado aproximado;
  \item número de puntos globalmente no dominados;
  \item distancia generacional inversa (IGD) contra el frente global combinado;
  \item spacing, como uniformidad de separación local entre puntos;
  \item dispersión o extensión del frente en el espacio \((u_1,u_2,u_3)\).
\end{itemize}

\subsection{Relajaciones PL y guía adaptativa del MMC-MO}

Para evitar una inicialización arbitraria, antes de iniciar la sociedad de
compositores se resuelven una sola vez, por instancia del problema, las
relajaciones lineales monoobjetivo asociadas a:
\[
u_1=1-f_1,\qquad u_2=f_2,\qquad u_3=1-f_3,
\]
y un punto balanceado. En la relajación se permite:
\[
0\leq x_v\leq 1,\qquad 0\leq r_a\leq 1.
\]
Las soluciones fraccionales se redondean seleccionando los nodos de mayor valor
\(x_v\), respetando el presupuesto \(k\). Estas soluciones no son el resultado
final del MMC-MO; son una \textbf{guía inicial} o memoria de referencia:
\[
M_0=\{C^{LP}_{u_1},C^{LP}_{u_2},C^{LP}_{u_3},C^{LP}_{bal}\}.
\]

Durante la composición, el MMC-MO no vuelve a resolver los PL. La memoria guía
se actualiza con soluciones no dominadas encontradas por los compositores:
\[
M_{t+1}=\operatorname{ND}(M_t\cup P_t),
\]
donde \(P_t\) es el conjunto de soluciones evaluadas en la iteración \(t\) y
\(\operatorname{ND}\) conserva soluciones factibles no dominadas bajo el criterio
de Pareto con manejo de infactibilidad. Esta estrategia hace que la guía inicial
derivada de PL oriente la búsqueda, pero que el método conserve su carácter
heurístico multiobjetivo.

\section{Salidas}

El sistema exporta:

\begin{itemize}
  \item eventos narrativos por documento;
  \item nodos de la red narrativa;
  \item aristas ponderadas;
  \item solución del cubridor nodal;
  \item frente de Pareto;
  \item hipervolumen;
  \item superficie empírica \(u_1\times u_3\to u_2\) para visualizar el frente
  tridimensional sin confundirlo con una proyección plana;
  \item métricas de cobertura, peso nodal, costo de aristas y valor objetivo de referencia.
\end{itemize}

\section{Limitaciones metodológicas}

El sistema depende de la calidad de los índices abiertos y de la posibilidad de
descargar texto. Crossref, OpenAlex, Google News RSS, GDELT y Reddit no
representan universos completos; representan mecanismos de muestreo auditables.
La extracción completa se condiciona por permisos de fuente: antes de descargar
HTML se consulta \texttt{robots.txt}; si la fuente es paywall, parcial o no
permite descarga, el registro se conserva sólo como metadato o señal
\texttt{ok\_partial}. La limpieza automática reduce ruido, pero no sustituye la
revisión manual de una muestra. Por ello, las conclusiones publicables deben
reportar consulta, variantes, años, región, fuentes activas, exclusiones, tasa
de errores de descarga, proporción de \texttt{ok}/\texttt{ok\_partial} y criterios
de validación humana.

\section{Reproducibilidad}

Todo el análisis corre localmente. Para ejecutar el solver:

\begin{verbatim}
python3 scripts/run_cover_solver.py \
  --input news_output \
  --method mmc \
  --max-nodes 12 \
  --edge-cost-weight 1.0
\end{verbatim}

\end{document}
